\documentclass[11pt, a4paper]{article}

% ─── PREAMBLE ───────────────────────────────────────────────────────────────
\usepackage[a4paper, top=2.5cm, bottom=2.5cm, left=2.5cm, right=2.5cm, headheight=15pt]{geometry}
\usepackage[utf8]{inputenc}
\usepackage[T1]{fontenc}
\usepackage[romanian]{babel}
\usepackage{xcolor}
\usepackage{fancyhdr}
\usepackage{booktabs}
\usepackage{tabularx}
\usepackage{longtable}
\usepackage{array}
\usepackage{enumitem}
\usepackage{parskip}
\usepackage[expansion=false]{microtype}
\usepackage[hidelinks]{hyperref}
\usepackage{multirow}
\usepackage{setspace}
\usepackage{titlesec}

\definecolor{vsnavy}{RGB}{15, 36, 71}
\definecolor{vsblue}{RGB}{37, 99, 235}
\definecolor{vsgray}{RGB}{100, 116, 139}
\definecolor{vsgreen}{RGB}{22, 163, 74}

\titleformat{\section}{\Large\bfseries\color{vsnavy}}{\thesection.}{0.5em}{}
\titleformat{\subsection}{\large\bfseries\color{vsnavy}}{\thesubsection.}{0.5em}{}
\titleformat{\subsubsection}{\normalsize\bfseries\color{vsnavy}}{\thesubsubsection.}{0.5em}{}
\titlespacing*{\section}{0pt}{1.6em}{0.8em}
\titlespacing*{\subsection}{0pt}{1.2em}{0.5em}
\titlespacing*{\subsubsection}{0pt}{1em}{0.4em}

\pagestyle{fancy}
\fancyhf{}
\fancyhead[L]{\small\color{vsgray}Caiet de practică — Vanadium Systems SRL}
\fancyhead[R]{\small\color{vsgray}Jipa Ștefan Ciprian}
\fancyfoot[C]{\small\color{vsgray}\thepage}
\renewcommand{\headrulewidth}{0.4pt}
\renewcommand{\footrulewidth}{0pt}

\hypersetup{
  pdftitle={Caiet de practică — Vanadium Systems SRL},
  pdfauthor={Jipa Ștefan Ciprian},
  pdfsubject={Practică de specialitate, Anul II, Grupa 4LF442}
}

\begin{document}

% ─── COPERTĂ ──────────────────────────────────────────────────────────────
\begin{titlepage}
\begin{center}

\vspace*{1.5cm}
{\large UNIVERSITATEA TRANSILVANIA DIN BRAȘOV}\\[0.2em]
{\large Facultatea de Inginerie Electrică și Știința Calculatoarelor}\\[0.2em]
{\normalsize Specializarea Automatică și Informatică Aplicată}

\vspace{2.5cm}
\rule{0.5\linewidth}{0.4pt}\\[0.6em]
{\Huge\bfseries\color{vsnavy} Caiet de practică}\\[0.6em]
{\Large Practică de specialitate}\\[0.6em]
\rule{0.5\linewidth}{0.4pt}

\vspace{2.5cm}
{\large
\begin{tabular}{rl}
Student: & Jipa Ștefan Ciprian \\[0.3em]
Grupa: & 4LF442 \\[0.3em]
An de studiu: & Anul II \\[0.3em]
Perioada practicii: & Iunie -- August 2026 \\[0.3em]
Firma / locul de practică: & Vanadium Systems SRL \\[0.3em]
Funcție în firmă: & Fondator \& Full-stack Developer \\
\end{tabular}
}

\vfill
{\small\color{vsgray}
Document elaborat în locul caietului de practică standard, cu acordul\\
prodecanului responsabil de practică, conform cerințelor comunicate:\\
tehnologii de dezvoltare utilizate, proiectele realizate și fluxul de lucru al unui proiect.}

\vspace{1cm}
{\small 2026}

\end{center}
\end{titlepage}

% ─── CUPRINS ────────────────────────────────────────────────────────────────
\setcounter{tocdepth}{2}
\tableofcontents
\newpage

% ─── INTRODUCERE ────────────────────────────────────────────────────────────

\section*{Context}
\addcontentsline{toc}{section}{Context}

Vanadium Systems SRL este o companie de software fondată și administrată de mine, care
dezvoltă produse SaaS verticale și orizontale pentru piața din România. Fiind fondator unic
(fără angajați), activitatea de practică a constat în dezvoltarea end-to-end a acestor produse:
discuții cu operatorii din domeniu, arhitectură, backend, frontend, infrastructură, testare,
securitate și mentenanță în producție: toate etapele unui produs software real, de la ideea
inițială până la clientul care îl folosește zilnic.

Acest document acoperă, conform cerinței comunicate de prodecan în locul caietului standard:
tehnologiile de dezvoltare folosite, proiectele realizate (fără nume de clienți) și fluxul de
lucru urmat pentru un proiect. Am extins fiecare secțiune cu suficient detaliu tehnic cât să
reflecte efectiv volumul de muncă din spatele ei, nu doar un rezumat de o linie pe subiect.

% ─── SECTIUNEA 1: TEHNOLOGII ─────────────────────────────────────────────────

\section{Tehnologii de dezvoltare folosite}

\subsection{Principii arhitecturale}

Înainte de stack-ul concret, merită explicate două decizii arhitecturale care se regăsesc
identic în toate platformele SaaS dezvoltate. Restul alegerilor tehnologice de mai jos sunt,
în mare parte, consecințe ale acestor două.

Prima este multi-tenancy cu izolare completă a datelor. O singură instanță de aplicație
(frontend și backend) deservește toți clienții, dar fiecare interogare către baza de date
filtrează explicit după identificatorul de client, \texttt{tenant\_id}. Fără excepții. E decizia
cu cel mai mare impact din tot codul: o singură scăpare aici înseamnă date dintr-o firmă vizibile
în alta, ceea ce e inacceptabil atât legal (GDPR) cât și contractual.

A doua este izolarea pe module de business. Fiecare capacitate funcțională (colectări, facturare,
rapoarte etc.) e un modul separat, cu propriul cod de rutare, logică și schemă de date. Modulele
nu importă logică unele din altele; când chiar trebuie să comunice, o fac explicit, prin baza de
date comună sau printr-o magistrală internă de evenimente, niciodată prin apel direct. Motivul e
practic: fiind un singur dezvoltator, am nevoie să pot modifica un modul fără să fiu obligat să
țin minte cum funcționează tot restul aplicației în același timp.

\subsection{Stack principal (aplicații SaaS)}

\renewcommand{\arraystretch}{1.25}
\begin{tabularx}{\textwidth}{>{\bfseries}l X}
\toprule
Nivel & Tehnologii \\
\midrule
Frontend & React 18, TypeScript, Vite, Tailwind CSS, shadcn/ui, TanStack Query, React Hook Form + Zod, Recharts, i18next \\
Backend & Node.js, Fastify, TypeScript, Drizzle ORM, PostgreSQL, node-cron, Puppeteer (generare PDF), MinIO (stocare fișiere), Nodemailer / Resend (e-mail) \\
Mobil & Capacitor (aceeași bază de cod React, aplicație Android/APK) \\
Infrastructură & Docker \& Docker Compose, Nginx, Certbot (SSL), Hetzner VPS, Cloudflare (DNS, Pages, Functions, Turnstile) \\
Testare \& CI/CD & Vitest, Playwright (E2E + accesibilitate), GitHub Actions (lint, typecheck, test, build, deploy) \\
\bottomrule
\end{tabularx}

\vspace{0.8em}
\noindent Câteva dintre aceste alegeri au o motivație destul de concretă, legată de faptul că
sunt un singur dezvoltator care răspunde și de operare, nu doar de scris cod, așa că le detaliez
puțin. Fastify e ales împreună cu Zod pe fiecare rută, astfel încât validarea cererii și a
răspunsului se face direct la nivel de framework, nu prin cod de validare scris manual pe care
îl pot uita undeva. Drizzle ORM, în loc de SQL brut, generează tipuri TypeScript direct din
schema bazei de date; dacă redenumesc o coloană, obțin erori de compilare imediat, în loc să
descopăr problema abia la runtime, în producție.

Pentru infrastructură am ales Docker Compose pe server propriu (Hetzner, Germania) în locul unei
platforme cloud gestionate, din două motive: cost fix lunar previzibil (contează pe un produs cu
marjă mică per client) și rezidența datelor în UE, cerută explicit de clienții din România pentru
conformitate GDPR. Iar pentru mobil folosesc Capacitor, ca să rulez aceeași bază de cod React și
ca aplicație Android. Întreținerea a două codebase-uri separate, web și mobil nativ, pur și
simplu nu e sustenabilă când lucrezi singur.

\vspace{0.5em}
\subsection{Automatizări \& produse orientate pe AI conversațional}

\begin{tabularx}{\textwidth}{>{\bfseries}l X}
\toprule
Categorie & Tehnologii \\
\midrule
Orchestrare fluxuri & n8n, Make \\
Voce AI & Vapi \\
Mesagerie & WhatsApp Business API, Gmail API, Google Calendar API \\
\bottomrule
\end{tabularx}

\vspace{0.5em}
\noindent Aceste produse nu sunt aplicații monolitice ca VanaGreen, ci fluxuri de automatizare
construite peste API-uri oficiale (nu integrări neoficiale, fragile) și orchestrate prin n8n sau
Make. Alegerea e motivată simplu, de viteză: un flux nou poate fi construit și pus în producție
în câteva zile, nu săptămâni. Contează, pentru că e vorba de un produs orizontal, vândut la mai
multe firme mici, fiecare cu cerințe ușor diferite de configurat.

\subsection{Instrumente de dezvoltare asistată de AI}

O parte semnificativă a fluxului de dezvoltare folosește un agent de programare AI
(\textbf{Claude Code}, Anthropic) ca instrument principal de scriere și refactorizare a codului,
sub supraveghere și revizuire umană directă la fiecare pas. Acest lucru este menționat separat
pentru că schimbă efectiv fluxul de lucru descris în secțiunea~4: planificarea și specificarea
cerințelor devin explicite și scrise (nu doar gândite), iar revizuirea de cod și verificarea
automată a calității devin pași obligatorii înaintea oricărei livrări, nu opționali.

Concret, agentul este folosit pentru:
\begin{itemize}[leftmargin=1.5em, itemsep=0.2em]
  \item implementarea de funcționalități noi pornind de la un plan scris în prealabil (structură DB, endpoint-uri backend, componente frontend);
  \item revizuire de cod automată (bug-uri, duplicare, simplificare) înainte de merge;
  \item scriere și menținere de teste automate;
  \item generare și actualizare de documentație tehnică internă (jurnal de decizii per director/modul).
\end{itemize}

Un plan scris (înainte de orice implementare) arată, redus la esențial, astfel:

\begingroup\small
\begin{verbatim}
Feature: prag de avertizare la colectari de la persoane fizice
Scop: evitarea depasirii plafonului legal de 1000 lei / persoana / an
      fara ca operatorul sa observe manual suma cumulata

Schema DB: adaugare coloana `cumulative_amount` (calculata) pe `persons`
Backend:   validare in service-ul de creare colectare, inaintea de commit
Frontend:  warning inline in formular, cu optiune de split automat
Cazuri de test: sub prag / exact la prag / peste prag / split ulterior
\end{verbatim}
\endgroup

\vspace{0.3em}
\noindent Fiecare linie din acest plan devine, la implementare, un pas verificabil separat --
diferența esențială față de a scrie cod direct, fără o etapă explicită de specificare.

\subsection{Securitate \& conformitate}

Standardul de securitate e identic pe toate platformele SaaS, indiferent de dimensiunea
clientului. La autentificare, JWT-ul stă exclusiv într-un cookie \texttt{HttpOnly}, niciodată în
\texttt{localStorage}; MFA prin TOTP e obligatoriu pentru rolurile cu acces la date sensibile, un
singur dispozitiv rămâne activ per utilizator, iar sesiunea expiră atât absolut cât și după
inactivitate. La nivel de rută avem CSRF prin double-submit pe toate mutațiile, rate limiting
(mai agresiv pe autentificare) și blocare temporară după încercări repetate eșuate de login.

Câmpurile cu date personale sunt criptate end-to-end: cheia se derivă cu PBKDF2 direct în
browser, datele ajung pe server deja criptate cu AES-GCM, iar backend-ul nu le vede niciodată în
clar. Toate acțiunile sensibile intră într-un jurnal de audit imuabil. La nivel de rețea, HTTPS
e impus (HSTS activ), iar baza de date și storage-ul intern (MinIO) nu sunt expuse public — se
accesează doar prin backend, cu autentificare. Backup-ul rulează automat, zilnic.

Secretele nu ajung niciodată în cod, în repository sau în variabile expuse către client; sunt
gestionate exclusiv prin secretele platformei de hosting, iar fiecare rulare de integrare
continuă scanează automat după secrete uitate accidental în cod. Pentru platformele cu date GDPR
relevante, orice modificare care atinge câmpuri de date personale, generarea documentelor legale
sau regulile de reținere/ștergere a datelor trece printr-o etapă de revizuire separată, dedicată
strict acestor aspecte (detaliată în secțiunea~4.5) — nu se amestecă cu revizuirea tehnică
obișnuită.

% ─── SECTIUNEA 2: ARHITECTURA ────────────────────────────────────────────────

\section{Arhitectura tehnică a unei platforme SaaS}

Ca să nu rămână abstract ce înseamnă ``modul izolat'' menționat mai sus, arăt aici, la nivel de
cod, cum sunt structurate concret platformele de tip VanaGreen sau VanaDash.

\subsection{Structura pe module}

Fiecare capacitate de business (colectări, facturare, rapoarte etc.) este un folder de sine
stătător, cu exact trei fișiere responsabile de roluri diferite:

\renewcommand{\arraystretch}{1.3}
\begin{tabularx}{\textwidth}{>{\bfseries}l X}
\toprule
Fișier & Responsabilitate \\
\midrule
\texttt{routes.ts} & Definește endpoint-urile HTTP, validarea Zod a cererii/răspunsului și
rolurile necesare pentru acces. Nu conține logică de business. \\
\texttt{service.ts} & Logica efectivă: interogări către baza de date, reguli de validare,
calcule. Este singurul loc care ``știe'' cum funcționează modulul respectiv. \\
\texttt{schema.ts} & Structura tabelelor din acest modul, definită cu Drizzle ORM (sursă unică
de adevăr, din care se generează atât tipurile TypeScript cât și migrațiile SQL). \\
\bottomrule
\end{tabularx}

\vspace{0.8em}
\noindent Un exemplu simplificat (fără date reale), reprezentativ pentru tiparul folosit
consecvent în toate modulele backend:

\begingroup\small
\begin{verbatim}
// routes.ts
fastify.patch('/:id', {
  preHandler: [authenticate, requireRole(['operator'])],
  schema: { body: UpdateBodySchema, response: { 200: ItemResponseSchema } },
}, async (req) => {
  return { success: true, data: await service.update(req.tenantId, req.params.id, req.body) }
})

// service.ts -- actualizare cu izolare per client (tenant)
export async function update(tenantId: string, id: string, input: UpdateInput) {
  const [row] = await db.update(table)
    .set({ ...input, updatedAt: new Date() })
    .where(and(eq(table.id, id), eq(table.tenantId, tenantId), isNull(table.deletedAt)))
    .returning()
  if (!row) throw createError({ statusCode: 404, code: 'NOT_FOUND' })
  return row
}
\end{verbatim}
\endgroup

\vspace{0.5em}
\noindent Observație cheie: \texttt{tenantId} apare explicit în fiecare interogare care atinge
date de business. Nu e opțional și nu e ``adăugat mai târziu''. E convenția cu cea mai mare
importanță din tot codul, pentru că o singură interogare fără acest filtru e suficientă pentru
o scurgere de date între clienți.

\subsection{Convenții de bază de date}

Fiecare tabel de business respectă aceeași structură minimă, indiferent de proiect:

\begingroup\small
\begin{verbatim}
id          uuid primary key default gen_random_uuid()
tenant_id   uuid not null references tenants(id)
created_at  timestamp not null default now()
updated_at  timestamp not null default now()
deleted_at  timestamp nullable   -- stergere logica (soft-delete), niciodata hard-delete
\end{verbatim}
\endgroup

\vspace{0.3em}
\noindent Ștergerea definitivă a datelor de business (\texttt{DELETE} SQL) nu se folosește
niciodată din cod. Orice ``ștergere'' setează doar \texttt{deleted\_at}, păstrând istoricul
recuperabil, necesar atât pentru audit intern cât și pentru anumite obligații legale de
păstrare a evidențelor. Fiecare cheie străină are o politică explicită (\texttt{restrict},
\texttt{cascade} sau \texttt{set null}), stabilită deliberat, nu implicit.

\subsection{Convenții de denumire}

Consecvența denumirilor contează disproporționat de mult într-un cod întreținut de o singură
persoană, pe termen lung. Reduce timpul de reorientare atunci când revin, peste luni de zile,
la un modul pe care nu l-am mai atins.

\renewcommand{\arraystretch}{1.25}
\begin{tabularx}{\textwidth}{>{\bfseries}l X}
\toprule
Element & Convenție \\
\midrule
Fișiere backend & \texttt{kebab-case} \\
Componente React & \texttt{PascalCase} \\
Hooks React & \texttt{use-kebab-case} \\
Variabile/funcții & \texttt{camelCase}, fără abrevieri neclare \\
Booleeni & prefix \texttt{is}/\texttt{has}/\texttt{can} (ex. \texttt{isLoading}) \\
Rute API & \texttt{kebab-case}, plural, fără verbe (ex. \texttt{/api/collections}) \\
Coloane DB & \texttt{snake\_case} \\
\bottomrule
\end{tabularx}

% ─── SECTIUNEA 3: PROIECTE ───────────────────────────────────────────────────

\section{Proiecte realizate}

\subsection{VanaGreen --- platformă SaaS pentru operatori de reciclare (proiect principal)}

Platformă SaaS multi-tenant care digitalizează complet activitatea unui operator de colectare
deșeuri din România: înregistrarea colectărilor, gestiunea stocului, generarea automată a
documentelor legale și rapoartele de conformitate. Este cel mai amplu și mai vechi produs al
firmei, cel pe care s-a validat inițial arhitectura descrisă în secțiunea~3 și reutilizată apoi
în celelalte proiecte.

\subsubsection{Module funcționale}

\begin{itemize}[leftmargin=1.5em, itemsep=0.35em]
  \item \textbf{Colectări.} Flux digital complet de la sursă (companie, persoană fizică,
  instituție, campanie) până la stocare, cu avertismente automate pentru flux mixt de deșeuri pe
  aceeași colectare, depășirea plafonului legal de 1.000 lei pentru persoane fizice și posibile
  duplicate. Se pot despărți sau uni colectări ulterior, cu istoric complet al modificării.
  \item \textbf{Expedieri.} Stocul pleacă spre colectori sau procesatori doar după validarea
  automată a cantității disponibile — nu se poate expedia mai mult decât există fizic în stoc —
  iar documentele aferente se generează automat.
  \item \textbf{Tratare.} Deșeurile compuse se descompun în fracții valorificabile pe baza unor
  rețete predefinite per client (un frigider, de exemplu, iese ca fier, aluminiu, plastic și
  agent frigorific), cu stocul actualizat automat la finalizare.
  \item \textbf{Documente legale și rapoarte de conformitate.} NIR, proces-verbal de intrare și
  ieșire, borderou de plată cu calcul automat al reținerilor legale, anexe pentru deșeuri
  periculoase, toate cu numerotare configurabilă per client. La acestea se adaugă exporturi gata
  de depus pentru evidența deșeurilor, raportări de mediu și declarații fiscale legate de
  reținerile la sursă.
  \item \textbf{Facturare.} Emitere, stornare, urmărire plăți, regim de TVA detectat automat și
  integrare directă cu e-Factura ANAF (SPV).
  \item \textbf{Containere și flotă.} Istoricul mișcărilor (ridicare, predare, relocare) pentru
  containerele proprii și cele închiriate, plus alerte de mentenanță pentru ITP, asigurări,
  revizii și verificări metrologice ale cântarelor.
  \item \textbf{Aplicație mobilă.} Aplicație Android pentru echipele de teren, cu funcționare
  offline și sincronizare automată la revenirea conexiunii: colectări, fotografii, program de
  ridicări, toate introduse direct din teren.
  \item \textbf{API public și webhook-uri}, pentru integrare cu sisteme ERP/WMS pe care clientul
  le are deja, prin REST documentat și evenimente în timp real.
  \item \textbf{White-label:} logo, culori și domeniu propriu configurabile per client, ca
  aplicația să poarte brandul firmei client, nu al Vanadium Systems.
\end{itemize}

\subsubsection{Securitate specifică domeniului}

Pe lângă măsurile generale descrise în secțiunea~2.5, VanaGreen tratează separat câmpurile de
date personale ale persoanelor fizice furnizoare de deșeuri (nume, CNP, date vehicul) --
criptate end-to-end, niciodată vizibile în clar pe server.

\subsubsection{Rezultate operaționale}

Cifre reale, măsurate din utilizarea în producție a platformei:

\renewcommand{\arraystretch}{1.3}
\begin{tabularx}{\textwidth}{X X}
\toprule
Indicator & Valoare \\
\midrule
Timp de la cântărire la document generat & sub 2 minute \\
Tipuri de documente legale generate automat & peste 15 \\
Digitalizare a fluxului operațional & 100\% (fără hârtie) \\
Conformitate & GDPR (criptare, minimizare, drept la ștergere) \\
\bottomrule
\end{tabularx}

\vspace{1.2em}
\subsection{VanaDash --- hub intern de monitorizare (proiect secundar)}

Dashboard intern, cu un singur utilizator (fără clienți externi, fără multi-tenancy), construit
pentru a agrega într-un singur loc datele din toate platformele Vanadium Systems, altfel
imposibil de urmărit central de o singură persoană care rulează în paralel mai multe produse.

\subsubsection{Ce face concret}

La bază stau niște conectori per platformă: module TypeScript dedicate, fiecare știind să
comunice cu API-ul unei singure platforme (de exemplu VanaGreen) și să extragă din ea metrici
standardizate. Peste ei rulează job-uri programate (\texttt{node-cron}) care fac sincronizarea
periodică și țin un jurnal de execuție pentru fiecare rulare: succes sau eșec, câte înregistrări,
ce eroare, dacă a fost cazul.

Datele adunate așa alimentează indicatorii financiari și operaționali centralizați, cu filtrare
pe perioade (luna curentă, ultimele 30 de zile, interval custom), plus un set de reguli de prag
configurabile care declanșează alertare automată — de exemplu, venit sub un anumit nivel
într-o zi. Singura excepție de la rolul strict de agregator este facturarea: facturile emise de
firmă sunt create și numerotate direct aici, apoi transmise ca e-Factură către ANAF.

\textbf{Tehnologii:} aceeași bază tehnică ca VanaGreen (React, Fastify, PostgreSQL, Drizzle ORM),
simplificată prin eliminarea completă a multi-tenancy-ului. E o decizie deliberată — o platformă
cu un singur utilizator nu are de ce să care după ea complexitatea izolării per client.

\vspace{1em}
\subsection{Alte produse SaaS live}

\subsubsection{Asistent virtual AI (recepționist automat)}

Preia apeluri telefonice, mesaje WhatsApp și e-mail simultan, non-stop, cu escaladare automată
la un operator uman pentru cazurile complexe. E gândit pentru firme mici care nu își permit o
recepție dedicată 24/7: clinici, cabinete, HoReCa, servicii B2B.

Vocea e personalizată pe tonul firmei și poate face programări direct în calendar, cu confirmare
pe SMS sau WhatsApp, plus răspunsuri la întrebări frecvente pe care clientul le configurează
singur, fără intervenție IT. Pe WhatsApp Business, confirmările și remindere-le automate merg
prin integrarea oficială a platformei, nu printr-un tool neoficial și fragil care se poate strica
oricând. Inbox-ul de e-mail e gestionat tot de AI: mesajele sunt clasificate automat (cereri de
ofertă, programări, reclamații, spam), iar un draft de răspuns e generat pe baza contextului
conversației. Când un caz depășește ce poate rezolva AI-ul, transferul către operatorul uman
vine cu tot contextul conversației atașat, ca să nu fie nevoie ca clientul să repete de la
capăt ce a spus deja.

\renewcommand{\arraystretch}{1.3}
\begin{tabularx}{\textwidth}{>{\bfseries}l X X}
\toprule
Plan & Include & Minute/lună \\
\midrule
Starter & Telefon SAU WhatsApp (un singur canal), programare automată & 200 \\
Business & Telefon + WhatsApp + E-mail, escaladare automată, jurnal complet & 600 \\
Enterprise & Canale nelimitate, integrare CRM/ERP existent, SLA contractual & negociat \\
\bottomrule
\end{tabularx}

\textbf{Tehnologii:} voce AI (Vapi), WhatsApp Business API, Gmail API, n8n, Make.

\subsubsection{Automatizări Gmail \& WhatsApp}

Automatizează procesele repetitive din inbox-ul unei firme: răspunsuri la cereri de ofertă,
urmărirea e-mailurilor rămase fără răspuns, sincronizare de calendar. E-mailurile primesc automat
etichetă și rutare pe categorii, cu un draft de răspuns sugerat, iar programările pe WhatsApp
primesc confirmare și remindere automate cu 24 de ore înainte. Google Calendar și CRM-ul clientului
sunt sincronizate bidirecțional, iar acolo unde firma are deja alte unelte în uz (Slack, Notion,
Sheets, Drive), le conectăm prin n8n, Make sau Zapier, ca să nu mai fie nevoie de copiat date
manual dintr-un loc în altul.

\textbf{Tehnologii:} Gmail API, WhatsApp Business API, Google Calendar API, n8n, Make, Zapier.

\subsubsection{Site-uri web}

Dezvoltare de site-uri de prezentare, landing page-uri și portaluri client pentru firme B2B,
orientate pe conversie și SEO tehnic, fără platforme de tip WordPress sau template-uri.

\renewcommand{\arraystretch}{1.3}
\begin{tabularx}{\textwidth}{>{\bfseries}l X l}
\toprule
Pachet & Ce include & Livrare \\
\midrule
Landing page & 1--3 pagini, un singur obiectiv de conversie, SEO de bază & 5 zile \\
Prezentare & până la 7 pagini, SEO tehnic complet, CMS pentru editare & 7--10 zile \\
Avansat & 7+ pagini, portal client cu autentificare, integrări API & la stabilit \\
\bottomrule
\end{tabularx}

\textbf{Tehnologii:} React, TypeScript, Vite, Tailwind CSS, Cloudflare Pages \& Functions.

\vspace{1em}
\subsection{Proiect în dezvoltare: platformă SaaS pentru service-uri auto}

În fază de concept și dezvoltare inițială, acest proiect reutilizează arhitectura deja validată
în producție la VanaGreen (aceeași structură de module, aceeași abordare multi-tenant), aplicată
unui domeniu diferit: service-uri auto, spălătorii, vulcanizări, tinichigerie-vopsitorie și
dezmembrări.

Obiectivul e să înlocuiască fluxul actual, bazat pe deviz trimis pe WhatsApp și fișă de lucru pe
hârtie, cu un flux digital complet. Doi diferențiatori pe care jucătorii existenți din piața
românească nu îi au încă îmi par cei mai importanți de urmărit. Primul e inspecția digitală a
vehiculului (DVI): un checklist pe tabletă pentru tehnician, cu cod de culoare pe severitate și
fotografii sau video atașate, unde clientul aprobă separat fiecare element din deviz, nu un
``da'' generic pentru tot. Al doilea sunt indicatori reali pentru proprietarul afacerii — KPI-uri
operaționale pe care service-urile auto din România pur și simplu nu le au azi disponibile
automat.

Un singur tenant poate activa simultan mai multe tipuri de flux, de exemplu service, spălătorie
și vulcanizare sub același acoperiș — situație frecventă pe piața locală, și unul dintre motivele
pentru care arhitectura modulară de la VanaGreen se potrivește direct aici.

% ─── SECTIUNEA 4: WORKFLOW ───────────────────────────────────────────────────

\section{Workflow-ul unui proiect}

Fluxul de lucru urmat pentru dezvoltarea și mentenanța fiecăruia dintre proiectele de mai sus
este structurat în etape clar delimitate, indiferent de mărimea funcționalității livrate.

\subsection{Discuția cu omul din spatele nevoii}

Orice proiect pornește de la o discuție, nu de la o presupunere. Înainte să scriu o singură
linie de cod sau să aleg vreo tehnologie, mă așez de vorbă, pe îndelete, cu omul care are nevoia
reală: operatorul care completează hârtii, proprietarul afacerii, persoana care simte zilnic
frustrarea unui proces greoi. Îl las pe el să-mi explice cum lucrează azi, unde pierde timp
sau bani, ce a încercat deja și nu a funcționat, ce și-ar dori de fapt.

Nu construiesc plecând de la ce cred eu că ar trebui să existe sau de la ce ``se poartă'' în
piață, ci de la nevoia reală a omului din fața mea. Discuția asta e primul pas, obligatoriu,
înaintea oricărei linii de plan tehnic. Tot ce urmează mai jos există ca să răspundă la ce am
aflat aici.

Un exemplu, fără date reale de client: un operator povestește că, la finalul fiecărei luni, un
angajat petrece 3--8 ore calculând manual, dintr-un registru pe hârtie, un raport pe care trebuie
să-l depună la autoritate. Din discuția asta reiese nevoia reală, care nu e ``vreau o aplicație'',
ci ``vreau să nu mai pierd o zi de muncă în fiecare lună pe un calcul repetitiv''. Abia acum are
sens să vorbim despre ce se construiește tehnic.

\subsection{Planificare}

Orice funcționalitate non-trivială pornește de la un plan scris explicit, înainte de a scrie
cod. Un plan conține, la minimum: ce anume se schimbă și de ce (pornind direct din discuția de
la pasul anterior), ce module sunt afectate, ce schimbări de schemă de date sunt necesare și ce
compromisuri arhitecturale există, ce se câștigă și ce se pierde. Planul trece printr-o etapă de
aprobare explicită înainte de a ajunge la implementare; nimic nu se implementează ``din mers'',
pe măsură ce apar idei noi în timpul scrisului de cod.

\subsection{Implementare pe etape (``wave'')}

Implementarea urmează o ordine fixă de dependențe, pentru a evita blocajele și integrările
premature:
\begin{enumerate}[leftmargin=2em, itemsep=0.3em]
  \item \textbf{Schema bazei de date și migrațiile.} Se stabilește întâi forma finală a datelor,
  pentru că o schimbare ulterioară de schemă costă mult mai mult decât una de cod aplicativ.
  \item \textbf{Backend.} Rute, logică de business și validare, modul cu modul, fiecare izolat,
  fără dependențe încrucișate între module (vezi secțiunea~3.1).
  \item \textbf{Frontend.} Componente și integrare cu un API deja funcțional și testat
  independent, ca un bug de UI să nu se confunde niciodată cu un bug de date.
  \item \textbf{Verificare finală și documentare,} ultimul pas înainte de testarea propriu-zisă.
\end{enumerate}

\subsection{Testare}

Testarea nu e un pas opțional la final, ci acoperă mai multe niveluri, fiecare cu alt scop.
Testele unitare și de integrare (Vitest) verifică logica de business izolat, în special calculele
și regulile de validare care n-au voie să greșească, cum ar fi calculul unui plafon legal.
Testele end-to-end, cu Playwright, simulează un flux complet de utilizator în browser, de la
login până la acțiunea finală, exact cum ar interacționa un client real cu aplicația.

Pe lângă acestea, verific automat și că interfața rămâne utilizabilă cu tastatura și cu cititoare
de ecran, nu doar vizual, și că fiecare pagină nouă respectă un prag minim de performanță
(Lighthouse), nu doar că ``merge''. Testarea manuală a fluxului complet rămâne ultimul filtru
înainte de livrare — mai ales pentru cazurile limită, greu de automatizat complet oricât ai
încerca.

\subsection{Revizuire de cod și securitate}

Fiecare schimbare trece printr-o etapă de revizuire de cod, urmărind trei categorii distincte:
corectitudine (există un bug?), simplificare/reutilizare (există cod duplicat sau o soluție mai
simplă?) și eficiență (există un cost de performanță nejustificat?).

Modificările care ating date personale, generarea documentelor legale, categorii reglementate
legal sau regulile de reținere/ștergere a datelor trec printr-o etapă suplimentară de revizuire,
dedicată exclusiv acestor aspecte, separată de revizuirea tehnică obișnuită, tocmai pentru că
o greșeală aici are consecințe legale, nu doar tehnice.

\subsection{Integrare continuă și livrare}

Fiecare modificare trece printr-un pipeline automat înainte de a putea fi livrată: verificare de
stil de cod (lint), verificare de tipuri, rulare a testelor automate, build complet, verificare
de vulnerabilități în dependențe și scanare automată de secrete accidental incluse în cod.
Nicio modificare nu ajunge în producție dacă vreunul dintre acești pași eșuează.

Livrarea propriu-zisă este automatizată și diferă în funcție de tipul de produs: platformele
SaaS rulează în containere Docker pe server propriu (Hetzner), cu reverse proxy Nginx și
certificate SSL reînnoite automat; site-urile de prezentare se publică automat, direct din Git,
către Cloudflare Pages, fără pas manual de deploy.

\subsection{Documentare}

La finalul fiecărei modificări semnificative actualizez un jurnal de decizii tehnice, cu patru
elemente fixe: ce exista înainte, ce s-a schimbat, care e starea curentă și ce compromisuri au
fost asumate. Ideea e ca orice decizie arhitecturală să poată fi reconstituită ulterior fără să
depindă de memoria mea. Contează mai ales când mai multe produse independente (VanaGreen,
VanaDash și celelalte) trebuie, ocazional, să comunice între ele: orice schimbare a acestui
contract de integrare se documentează explicit, în ambele proiecte afectate, nu doar în cel
unde a fost făcută schimbarea.

% ─── CONCLUZII ────────────────────────────────────────────────────────────────

\section*{Concluzii}
\addcontentsline{toc}{section}{Concluzii}
\enlargethispage{6\baselineskip}

Cel mai mult mi-a rămas din perioada asta faptul că nu pot da vina pe altcineva pentru o decizie
proastă. Când ești singurul dezvoltator, cel care vorbește cu omul din spatele nevoii e aceeași
persoană care proiectează schema bazei de date, scrie backend-ul, scrie frontend-ul, testează,
verifică securitatea și repară ce se strică duminică seara. Un compromis arhitectural asumat azi
devine, garantat, problema mea de peste trei luni. Curricula facultății nu prea are cum să simuleze
genul ăsta de presiune — și cred că e exact partea care lipsește cel mai des dintr-un proiect de
laborator, oricât de bine ar fi construit tehnic.

Nu susțin că modul ăsta de lucru e ideal sau că merge scalat la nesfârșit; sunt lucruri pe care
le-aș fi făcut altfel dacă aș fi avut de la cine să cer o a doua opinie tehnică înainte de a lua o
decizie. Dar practica desfășurată în Iunie -- August 2026 a fost, cel puțin, reală: produse care
rulează în producție, cu clienți activi, nu un exercițiu de laborator scris ca să bifeze o
cerință. Proiectul aflat acum în dezvoltare, platforma pentru service-uri auto, continuă pe
același drum.

\vspace{1em}
\noindent\rule{\linewidth}{0.4pt}
\vspace{0.5em}

\noindent\small\color{vsgray}
Document întocmit de Jipa Ștefan Ciprian, Vanadium Systems SRL, în cadrul practicii de
specialitate desfășurate în perioada Iunie -- August 2026.

\end{document}
